% !TEX root = ../../main.tex

\section{Projected Wald Example in Matrix Form}
\label{app:projected-wald-example}

This appendix expands the worked sibling example from the Method section in
full matrix form. Let
\begin{equation}
z =
\begin{bmatrix}
2.0\\
-1.0\\
0.5\\
1.0
\end{bmatrix},
\qquad
P =
\begin{bmatrix}
1&0&0&0\\
0&1&0&0\\
0&0&1&0
\end{bmatrix}.
\end{equation}
The projected vector is
\begin{equation}
y = Pz =
\begin{bmatrix}
1&0&0&0\\
0&1&0&0\\
0&0&1&0
\end{bmatrix}
\begin{bmatrix}
2.0\\
-1.0\\
0.5\\
1.0
\end{bmatrix}
=
\begin{bmatrix}
2.0\\
-1.0\\
0.5
\end{bmatrix}.
\end{equation}
The plain quadratic form is therefore
\begin{equation}
y^\top y
=
\begin{bmatrix}
2.0 & -1.0 & 0.5
\end{bmatrix}
\begin{bmatrix}
2.0\\
-1.0\\
0.5
\end{bmatrix}
=
5.25.
\end{equation}
If the projected coordinates were all unit-variance Gaussian noise, this would
be compared directly to $\chi^2_3$.

Now suppose the first two coordinates are PCA-derived and have eigenvalue
matrix
\begin{equation}
\Lambda =
\begin{bmatrix}
4&0\\
0&1
\end{bmatrix},
\qquad
\Lambda^{-1} =
\begin{bmatrix}
1/4&0\\
0&1
\end{bmatrix}.
\end{equation}
Writing
\begin{equation}
y_{\mathrm{pca}} =
\begin{bmatrix}
2.0\\
-1.0
\end{bmatrix},
\qquad
y_{\mathrm{pad}} =
\begin{bmatrix}
0.5
\end{bmatrix},
\end{equation}
the whitened statistic is
\begin{equation}
y_{\mathrm{pca}}^\top \Lambda^{-1} y_{\mathrm{pca}}
+
y_{\mathrm{pad}}^\top y_{\mathrm{pad}}
=
\begin{bmatrix}
2.0 & -1.0
\end{bmatrix}
\begin{bmatrix}
1/4&0\\
0&1
\end{bmatrix}
\begin{bmatrix}
2.0\\
-1.0
\end{bmatrix}
+
0.25
=
2.25.
\end{equation}
This is the per-component whitening path: the PCA coordinates are each divided
by their own null variance before being squared and summed.

For the orthonormal projected-Wald calibration, the same raw statistic $5.25$
is kept. The PCA eigenvalues select the directions, but the null law is not
weighted by them because the projected coordinates are inner products with
orthonormal vectors. With three projected coordinates,
\begin{equation}
p = \Pr\!\left(\chi^2_3 \ge 5.25\right).
\end{equation}
The practical distinction is the same as in the main text: whitening changes
the coordinates before summation, whereas the implemented projected-Wald test
keeps the raw orthonormal projection and uses its chi-square dimension.
